\begin{afterword}
\summaryhead

The post describes Solomon Asch's experiments of the 1950s. A subject judges which of three lines matches a standard line, after a row of confederates has given the same wrong answer. Three quarters of Asch's subjects conformed at least once, and a third conformed more than half the time.

The author then asks whether conforming was irrational. By Aumann's agreement theorem, other people's beliefs are evidence, and an honest observer facing three others who all see the same diagram might reasonably give the majority more than even odds. The post argues that the details of later results rule out this reading. Conformity stops rising after three confederates; one dissenter cuts it sharply, and it returns if the dissenter switches sides; it is higher among women and among ingroup members; it is lower when the error is blatant and when answers are private. Subjects also deny that a dissenter influenced them. The pattern points to nervousness about being the odd one out.

\responsehead

The post asks a good question, whether going along with a unanimous group can be rational, and reports Asch's main figures correctly. Its answer depends on a list of later findings, which it reports less carefully than the figures.

The opening numbers give only part of Asch's result. Three quarters of the subjects conformed at least once, and a third more than half the time. Per answer, subjects went along a little over a third of the time, and Asch wrote that the majority effect was not ``even the strongest force at work.'' Most of the answers he recorded were correct.

The theorem is stated without one of its two conditions. Aumann's result holds for people who start from the same prior beliefs; people who start from different ones can agree to disagree.

The post tests the findings against two explanations, conformity as the rational use of evidence and conformity as a way to avoid sticking out, and its ``(all)'' allows that both are at work. That mixture is the standard account. Deutsch and Gerard named the two influences ``informational'' and ``normative'' in 1955, in a study that included an anonymous version of Asch's task. The post does not mention them.

Some findings are reported more simply than Asch reported them. His dissenter figures are shares of answers, and they depend on the kind of dissenter; a dissenter whose wrong answer was less extreme than the majority's cut the majority's effect by only about a third. The return of conformity when the dissenter defected was, in Asch's account, an effect of desertion: a partner who simply left had an effect that ``outlasted his presence.'' Several findings attributed to the meta-analysis were first reported by Asch. The meta-analysis supports the direction of the others.

In short: a good question and Asch's figures reported correctly, with a theorem stated without its condition that both parties start from the same beliefs, some findings simplified or credited to the wrong source, and the two motives it weighs left without the names Deutsch and Gerard gave them in 1955.
\end{afterword}
